% Options for packages loaded elsewhere
\PassOptionsToPackage{unicode}{hyperref}
\PassOptionsToPackage{hyphens}{url}
%
\documentclass[
]{article}
\usepackage{amsmath,amssymb}
\usepackage{lmodern}
\usepackage{iftex}
\ifPDFTeX
  \usepackage[T1]{fontenc}
  \usepackage[utf8]{inputenc}
  \usepackage{textcomp} % provide euro and other symbols
\else % if luatex or xetex
  \usepackage{unicode-math}
  \defaultfontfeatures{Scale=MatchLowercase}
  \defaultfontfeatures[\rmfamily]{Ligatures=TeX,Scale=1}
\fi
% Use upquote if available, for straight quotes in verbatim environments
\IfFileExists{upquote.sty}{\usepackage{upquote}}{}
\IfFileExists{microtype.sty}{% use microtype if available
  \usepackage[]{microtype}
  \UseMicrotypeSet[protrusion]{basicmath} % disable protrusion for tt fonts
}{}
\makeatletter
\@ifundefined{KOMAClassName}{% if non-KOMA class
  \IfFileExists{parskip.sty}{%
    \usepackage{parskip}
  }{% else
    \setlength{\parindent}{0pt}
    \setlength{\parskip}{6pt plus 2pt minus 1pt}}
}{% if KOMA class
  \KOMAoptions{parskip=half}}
\makeatother
\usepackage{xcolor}
\setlength{\emergencystretch}{3em} % prevent overfull lines
\providecommand{\tightlist}{%
  \setlength{\itemsep}{0pt}\setlength{\parskip}{0pt}}
\setcounter{secnumdepth}{-\maxdimen} % remove section numbering
\ifLuaTeX
  \usepackage{selnolig}  % disable illegal ligatures
\fi
\IfFileExists{bookmark.sty}{\usepackage{bookmark}}{\usepackage{hyperref}}
\IfFileExists{xurl.sty}{\usepackage{xurl}}{} % add URL line breaks if available
\urlstyle{same} % disable monospaced font for URLs
\hypersetup{
  hidelinks,
  pdfcreator={LaTeX via pandoc}}

\author{Dietmar Gerald Schrausser}
\date{09.05.2026}


\begin{document}

\hypertarget{foliumblat-vir}{%
\section{Foliu{[}m{]}/Blat VIr}\label{foliumblat-vir}}

Eine zeitgenössische deutsche Übersetzung wird präsentiert, die auf
einem Vergleich der bestehenden lateinischen, frühneuhochdeutschen
(ENHG) und englischen (Übersetzungs-) Texte basiert, um ein besseres
Verständnis zu ermöglichen. Insbesondere da der ENHG-Text im Vergleich
zum modernen Hochdeutschen schwer verständlich ist, da dieser fundierte
Deutschkenntnisse (als Muttersprache) sowie Kenntnisse verschiedener
Dialekte voraussetzt.

\hypertarget{bestimmtes-ist-uxfcber-die-planetensphuxe4re-und-die-umlaufbahnen-zu-erluxe4utern.}{%
\subsection{Bestimmtes ist über die Planetensphäre und die Umlaufbahnen
zu
erläutern.}\label{bestimmtes-ist-uxfcber-die-planetensphuxe4re-und-die-umlaufbahnen-zu-erluxe4utern.}}

\hypertarget{der-unterschied-zwischen-himmlischen-und-elementaren-umlaufbahnen1}{%
\subsubsection[Der Unterschied zwischen himmlischen und elementaren
Umlaufbahnen]{\texorpdfstring{Der Unterschied zwischen himmlischen und
elementaren
Umlaufbahnen\footnote{`orbitu{[}m{]}' und `umbkreys'.}}{Der Unterschied zwischen himmlischen und elementaren Umlaufbahnen}}\label{der-unterschied-zwischen-himmlischen-und-elementaren-umlaufbahnen1}}

Der gesamte Mechanismus des \emph{Korpus}\footnote{`COrporalis' und
  `leiplich'.} der Welt besteht aus \emph{zwei} Dingen: nämlich der (I)
\emph{himmlischen} Natur und der (II) \emph{elementaren} Natur. Darüber
hinaus ist der Himmel (I) in \emph{drei} \emph{Haupthimmel} unterteilt,
nämlich in den (i) \emph{feurigen}\footnote{`empyreum' und `feürigen'.},
den (ii) \emph{kristallinen} und in das (iii) \emph{Firmament}.

Innerhalb des Firmaments (iii), welches \emph{in der Tat}\footnote{`vero'.}
der Himmel der Sterne\footnote{`celu{[}m{]} stellarum' und `gestirnt
  himel'.} ist, befinden sich die sieben Umlaufbahnen der sieben
Planeten, das sind Saturn, Jupiter, Mars, Sonne, Venus, Merkur, Mond.
Unter dem Namen des kristallinen oder wässrigen\footnote{`aquei', in der
  mittelalterlichen Vorstellung handelte es sich dabei um die Sphäre
  über den Sternen, die oft mit den in der Genesis erwähnten
  \emph{Wassern über dem Firmament} in Verbindung gebracht wurde.}
Himmels (ii) versteht man jedoch den ersten Teil der
Urmaterie\footnote{`materie prime' und `erste{[}n{]} materi'}, der nach
Ansicht des Philosophen\footnote{`s{[}ecundum{]} philosophu{[}m{]}' und
  `nach sag des weysen' bezieht sich in diesem Kontext fast immer auf
  Aristoteles.} in zwei Umlaufbahnen gegliedert ist. Die obere
Umlaufbahn ist das \emph{primum mobile}. Die Natur dieser Umlaufbahn
ist, dass sie alle bewegt werden\footnote{`mouent{[}ur{]}' und `bewegt
  werden'.}, mit Ausnahme des feurigen Himmels (i), welcher ruht.

Tatsächlich\footnote{`vero' und `Aber'.} lässt sich die
\emph{elementare} Natur (II) in \emph{vier} Hauptsphären unterteilen,
nämlich (i) Feuer, (ii) Luft, (iii) Erde und (iv) Wasser. Die Sphäre des
Feuers (i) hat drei Abschnitte, den (a) \emph{obersten}, der als
\emph{feurig} bezeichnet wird, sowie den \emph{mittleren} und den
\emph{untersten}, der als olympisch\footnote{`olimpium' und `liecht'.}
bezeichnet wird.

In ähnlicher Weise hat die Luft (ii) drei Abschnitte, den (a)
\emph{obersten}, der als \emph{ätherisch}\footnote{`ethereum' und
  `scheynlich'.} bezeichnet wird und den (b) \emph{mittleren} sowie den
(c) \emph{untersten}, welche/r als \emph{luftig}\footnote{`aereum' und
  `lüftig'.} bezeichnet wird/werden\footnote{`dicit{[}ur{]}', singular
  vs.~`der mittel vn{[}d{]} underst lüftig'.}. Im obersten Bereich (a)
gibt es aufgrund der Nähe der Sonne Wärme und Licht, ähnlich verhält es
sich auch im untersten Teil (c), dies aber aufgrund der
\emph{Reflexion}\footnote{`rep{[}er{]}cussionem' und `wider scheyns'.}
der \emph{Strahlen von der Erde}\footnote{`ratio{[}rum{]} \ldots{} a
  terra' und `glentz von der erden'.}. Im mittleren Abschnitt (b)
jedoch, den die \emph{reflektierten Strahlen}\footnote{`rep{[}er{]}cussio
  radio{[}rum{]}' und `widerscheyn d{[}er{]} glentz'.} nicht erreichen
können, herrscht Kälte und Dunkelheit. Es wird gesagt, dass dieser von
\emph{Dämonen}\footnote{`demones' und `teüfel'.} bewohnt wird, welche in
diese \emph{dunkle}\footnote{`caliginosum' und `tunckeln'.} \emph{Luft}
verstossen\footnote{`detrusi' und `verstössen'.} sind.\footnote{findet
  sich insbesondere im Handbuch der christlichen Lehre, dem
  `Elucidarium', ein häufiger theologischer Bezug auf das
  \emph{atmosphärische Gefängnis}, in dem gefallene Engel bis zum
  Jüngsten Gericht verweilen, vgl. Epheser oder 2 Petrus 2,4 `da Gott
  die Engel, die gesündigt hatten, nicht verschonte, sondern sie in die
  \emph{Hölle} stieß und sie in \emph{finstere Ketten} legte,
  {[}\ldots{]}'. Man Beachten hier besonders die `\emph{finsteren
  Ketten}', welche die vollständige Gefangenschaft und \emph{Trennung
  von} \emph{Gottes Licht} symbolisieren.} Dort entstehen auch
Unwetter\footnote{`tempestates' und `ungestümigkeit'.} wie Donner,
Hagel, Schnee und ähnliches.\footnote{vgl. Francesco Maria Guazzo
  (\href{https://archive.org/details/compendiummalefi00guaz/page/n2/mode/1up}{1608}).}

Aus diesen bilde man \emph{zwölf Umlaufbahn}\footnote{`orbes' und
  `umbkrais'.} der Welt, die \emph{Erde und Wasser} (iii, iv) umgeben,
welche alle \emph{Himmel} genannt werden können. Doch diese
übertrifft\footnote{`excellit' und `ubertrifft'.} der \emph{Himmel der
Dreifaltigkeit}, \emph{Gott selbst}\footnote{`ipse deus' und `der got'.},
der \emph{unter allen} Dinge\footnote{`in omnibus' und `in allen'.} und
\emph{jenseits aller} Dinge\footnote{`super om{[}n{]}ia' und `über
  alle'.} ist.

Die Entfernungen der \emph{vorgenannten}\footnote{`predictorum' und
  `vorgenante{[}n{]}'.} Umlaufbahnen und Planeten betragen von der Erde
zum Mond 15625 Meilen; \emph{Meilen, das sind} \textbf{\emph{126}}
\emph{Stadien}\footnote{`miliaria. hec sunt stadia .cxxvi.', ergibt
  15625 Meilen × 126 = 1968750 Stadien, das ergibt bei \emph{attischen}
  Stadien, 0.18498 × 1968750 = 364179.375 km Entfernung von der Erde zum
  Mond. Bei aktueller \emph{maximaler} Entfernung (\emph{Apogäum}) von
  406740 km also eine Differenz von - 42560.625 km, bei \emph{minimaler}
  Entfernung (\emph{Perigäum}) von 356410 km eine Differenz von +
  7769.375 km. Dies \emph{entspricht} jedoch der Entfernung des Mondes vor etwa
  \emph{800 Millionen} Jahren (\emph{Tonium im Neoproterozoikum}, 363000
  bis 367000 km).}. Vom Mond zum Merkur sind es 7812.5 Meilen. Vom
Merkur zur Venus ist es die gleiche Entfernung. Von der Venus zur Sonne
sind es 22436 Meilen. Von der Sonne zum Mars sind es 15625 Meilen. Vom
Mars zum Jupiter sind es 7812 Meilen. Vom Jupiter zum Saturn ist es die
gleiche Entfernung. Vom Saturn zum Firmament sind es 22436 Meilen.
Daraus ergibt sich, dass die Entfernung von der Erde zum Sternenhimmel
109375 Meilen beträgt.

\hypertarget{die-unterscheidung-der-himmlischen-hierarchien-macht-oder-fuxfcrstentuxfcmer}{%
\subsubsection{Die Unterscheidung der himmlischen Hierarchien {[}Macht
oder
Fürstentümer{]}}\label{die-unterscheidung-der-himmlischen-hierarchien-macht-oder-fuxfcrstentuxfcmer}}

Was die \emph{himmlische} Natur (I) betrifft, so haben etliche Autoren
in der Tat eine \emph{dreifache} Unterscheidung getroffen, in eine (i)
\emph{überhimmlische} (oder das \emph{Übernatürliche}), eine (ii)
\emph{himmlische} und eine (iii) \emph{unterhimmlische}.

Die überhimmlische Hierarchie (i) existiere in drei Personen, wie manche
behauptet haben, \emph{jedoch fälschlicherweise}, denn laut
\textbf{Dionysios} impliziert \emph{Hierarchie} eine Ordnung, jedoch
eine solche Ordnung existiert in den drei Personen nicht absolut, außer
als Naturordnung.\footnote{`\ldots{} vt q{[}ui{]}dam dixerunt {[}et{]}
  male \ldots{}', von Pseudo-Dionysius Areopagita
  (\href{https://digital.bodleian.ox.ac.uk/objects/ee19a692-5066-4c84-95db-d93c4fc82b97/}{1350})
  definiert, ob der Begriff Hierarchie auf die Heilige Dreifaltigkeit
  angewendet werden kann.}

Die \emph{himmlische} (ii) ist in den \emph{Orden der Engel}, die
\emph{unterhimmlische} (iii) in \emph{heiligen Menschen}. Des Weiteren
ist die \emph{himmlische} Hierarchie (ii) in die (a) obere, die (b)
mittlere und die (c) untere Ebene unterteilt. Die obere Ebene (a)
umfasst drei Ordnungen, die (a1) \emph{Seraphim}, die (a2)
\emph{Cherubim} und die (a3) \emph{Thronengel}. Die ersten (a1)
erwägen\footnote{`{[}con{]}siderant' und `betrachten'.} Gottes
\emph{Wohlwollen}\footnote{`bonitate{[}m{]}' und `guttheit'.}, die
zweiten (a2) seine \emph{Größe}\footnote{`virtute{[}m{]}' und `kraft'.},
die dritten (a3) seine \emph{Angemessenheit}\footnote{`equitate{[}m{]}'
  und `gleicheit'.}. Dem gemäß \emph{liebt} Gott in den ersten (a1)
durch die \emph{Nächstenliebe}\footnote{`charitas' und `lib'.}, in den
zweiten (a2) erkennt er durch \emph{Wahrheit}, in den dritten (a3)
residiert\footnote{`sedet' und `sitzt'.} er in
\emph{Gerechtigkeit}\footnote{`equitas' und `gleicheit'.}.\footnote{von
  Pseudo-Dionysius Areopagita
  (\href{https://digital.bodleian.ox.ac.uk/objects/ee19a692-5066-4c84-95db-d93c4fc82b97/}{1350}),
  oft auch dem heiligen Gregor dem Großen zugeschrieben (vgl. Étaix,
  \href{https://books.google.com/books?id=r4ERnwEACAAJ}{1999}), auch
  adaptiert von Meister Eckhart
  (\href{http://www.eckhart.de/opus.htm}{1936}), widerspiegelt seine
  trinitarische Theologie, in welcher er die verschiedene Arten der
  \emph{Gotteswahrnehmung} mit den drei \emph{Personen} der
  Dreifaltigkeit in Einklang bringt.}

Die mittlere Hierarchie (b) umfasst (b1) \emph{Herrsch-}\footnote{`dominat{[}i{]}ones'
  und `herschengel'.}, (b2) \emph{Fürsten-}\footnote{`principatus' und
  `fürstengel'.} und (b3) \emph{Gewaltengel}\footnote{`potestates' und
  `gewaltengel', besser \emph{Fähigkeit}, \emph{Können}.}. Die ersten
(b1) regieren die Aufgaben der Engel, die \emph{folgenden}\footnote{`Sequentes'
  und `die andern'.} sind \emph{Berater}\footnote{`capitibunt
  p{[}re{]}sunt' und `pflegen der öbern'.} der Oberhäupter der Völker.
Die \emph{letzten} (b3) unterdrücken\footnote{`coercent' und
  `zwinge{[}n{]}'.} die \emph{Fähigkeit}\footnote{`potestatem' und
  `macht'.} der \emph{`Teufel'}\footnote{`demonu{[}m{]}' und `teüfel',
  \emph{der Verbrecher}.}. Ebenso \emph{herrscht} der Herr unter den
ersten (b1) als \emph{Majestät}, unter den folgenden\footnote{`in
  sequentibus' und `in den andern'.} (b2) \emph{regiert} er durch
\emph{ein Fürstentum}\footnote{`principat{[}us{]}' und `fürstenthumb'.},
durch die letzten (b3) wird er \emph{als `gesund' gehalten}\footnote{`tenet{[}ur{]}
  vt salus' und `gehalte{[}n{]} als das hail'.}.

Die untere Hierarchie (c) umfasst ebenfalls drei Ordnungen, (c1)
\emph{Kraftengel}\footnote{`virtutes' und `kreftengel'.}, (c2)
\emph{Erzengel} und (c3) \emph{Engel}. Den Kraftengeln (c1) obliegt die
Vollbringung großer Wunder, den Erzengeln (c2) die \emph{Verkündigung}
wichtiger Angelegenheiten, die Engel (c3) sorgen für den \emph{Schutz}
der Menschen. So wirkt Gott in den ersten (c1) als eine
\emph{Kraft}\footnote{`virt{[}us{]}' und `kraft'.}, in den
\emph{folgenden} (c2) offenbart er sich als \emph{Licht}, in den dritten
(c3) \emph{nährt} er als \emph{Inspiration}\footnote{`inspira{[}n{]}s'
  und `eyngeystender'.}.

Dies laut {[}Papst{]} \textbf{Gregor}.\footnote{`Hec dicta sunt
  s{[}ecundu{]}m Gregoriu{[}m{]}', in mittelalterlichen Handschriften
  üblich, kennzeichnete die vorhergehende Passage als zitiert vom Werk
  Papst Gregors I.} Nach \textbf{Dionysios}, jedoch bilden die
\emph{Kraftengel} (c1) die mittlere Ordnung der zweiten Hierarchie (b2),
die \emph{Fürstenengel} (b2) hingegen die erste Ordnung der dritten
Hierarchie (c1). Dies ist zu beachten, damit wird die \emph{Trinität der
göttlichen Person}\footnote{`trinitas p{[}er{]}sonarum diuinarum' und
  `trinitet der gottliche{[}n{]} person'.} in \emph{jeder} der
\emph{drei} genannten Hierarchien (i, ii, {[}iii{]}), auf allen Ebenen,
der oberen (1), mittleren (2) und unteren (3) gleichermaßen, so
wie\footnote{`vt patet'.} in der himmlischen Hierarchie (ii),
sichtbar.\footnote{Thomas von Aquin
  (\href{https://doi.org/10.3931/e-rara-15170}{1485}) führte dies in
  Werken wie der `Summa Theologica' weiter aus und folgte der
  dionysischen Struktur. Dass jede Hierarchie die \emph{Trinität
  göttlicher Personen} widerspiegelt, ist ein zentrales Thema der
  mittelalterlichen Scholastik.}

\hypertarget{von-der-zeit-oder-den-zeitaltern}{%
\subsubsection{Von der Zeit oder den
Zeitaltern}\label{von-der-zeit-oder-den-zeitaltern}}

Die Zeitalter der Welt werden sinnbildlich nach den
\emph{Altersabschnitten} des Menschen verstanden. Von den sechs
Zeitalter\footnote{vom Heiligen Augustinus
  (\href{https://doi.org/10.3931/e-rara-12425}{1489}) populär gemachtes
  christliches Geschichtsverständnis. In `De civitate Dei', legte er die
  Grundlage für diesen chronologischen Rahmen, während des Niederganges
  des West Römischen Imperiums, als eine umfassende christliche
  Interpretation von \emph{Menschheitsgeschichte} und dem Verhältnis
  zwischen der \emph{geistlichen und weltlichen} Welt.} der Welt,
beginnt das erste (I) mit der \emph{Welt Schöpfung} und dauert bis zur
Sintflut. Gemäß der hebräischen Darstellung\footnote{`hebraicam
  veritate{[}m{]}' und `nach hebreyscher warheit', bezog sich auf den
  \emph{ursprünglichen} hebräischen Text der Bibel, im Gegensatz zur
  griechischen \emph{Septuaginta}.} waren es 1656 Jahre\footnote{`.i656.'
  vs.~`.im.vc. lvi.', 1656 vs.~1556.}, nach der \textbf{Septuaginta}
(wie \textbf{Isidor} erklärt) und vielen anderen, denen wir hier folgen,
waren es jedoch 2242 Jahre.

Und so unterscheiden sie sich also um die 586 {[}\emph{686}{]} Jahre,
welche die Hebräer in diesem Zeitalter weniger haben.\footnote{charakteristisch
  für Isidor von Sevilla
  (\href{https://books.google.com/books?id=wb8Z_vlSyJwC}{1802}) oder
  Beda
  (\href{http://archivesetmanuscrits.bnf.fr/ark:/12148/cc76825p}{1125}),
  diese brachten die unterschiedlichen Zeitabläufe in Einklang.} Nach
dieser Berechnung starb \emph{Methusalem}\footnote{`matusale mortuus'
  und `matusaleh gestorben', ein häufiges chronologisches Problem in der
  mittelalterlichen Forschung, ob Methusalem die Sintflut überlebt
  hatte.} zwar vor der Sintflut, jedoch noch im selben Jahr, in dem die
Sintflut stattfand.

Das zweite\footnote{`S{[}e{]}c{[}un{]}da' und `Das andre'.} Zeitalter
(II) beginnt mit der Sintflut und dauert bis zur Geburt Abrahams, nach
den Hebräern dauerte es 292 Jahre\footnote{`292' vs.~`.ijc.lxxxij.', 292
  vs. 282.}, nach der Septuaginta 942 Jahre. Somit beträgt der
Unterschied 650 {[}\emph{660}{]} Jahre, welche die Hebräer wiederum
weniger haben. Doch worin der Grund für diese großen Unterschiede liegt,
\emph{konnte ich nicht herausfinden}\footnote{`inuenire no{[}n{]} potui'
  und `nicht möge{[}n{]} finden'.}.

Das Dritte Zeitalter (III) beginnt mit der Geburt Abrahams und dauerte
bis zum Beginn der Herrschaft Davids. Nach den Hebräern umfasste es 941
Jahre, nach der Septuaginta 940 Jahre. Das vierte Zeitalter (IV) beginnt
mit der Herrschaft Davids und dauert bis zur babylonischen
Gefangenschaft\footnote{`transmigrat{[}i{]}o{[}n{]}is babilonis' und
  `übergang babilonis'.}, es umfasst 484 Jahre nach den Hebräern und 485
Jahre nach der Septuaginta.

Das fünfte Zeitalter (V) beginnt mit der Babylonischen Gefangenschaft
genauer gesagt mit der Zerstörung Jerusalems und der Zerstörung des
darin befindlichen Tempels und dauerte bis zur Geburt Christi. Nach dem
oben genannten Modus, dem wir hier folgen, umfasst es 590 Jahre. Hier
herrscht eine große Uneinigkeit über die Berechnung der Jahre dieses
Zeitalters, manche berechnen sie in unterschiedlichen Weisen.

Das sechste Zeitalter (VI) beginnt mit der Geburt Christi und dauert bis
zum \emph{Ende der Welt}, dessen \emph{Termin}\footnote{`terminum' und
  `zil'.} allein Gott kennt. Dies wird \emph{Alter}\footnote{`senectus'
  und `das alt'.} oder die \emph{jüngste Stunde}\footnote{`hora
  nouissima' und `letzt stu{[}n{]}d', bekannt als erster Satz des `De
  Contemptu Mundi' von Bernhard von Cluny
  (\href{https://books.google.com/books?id=pAM_lfiF7EUC}{1864}), in dem
  er die Verderbnis der Welt und das Herannahen des Endes beschreibt.}
genannt. Zu diesen Zeitaltern aber mag man noch ein \emph{7. Zeitalter}
hinzufügen, jenes derer, die nun \emph{ruhen}\footnote{`quiescentium'
  und `die nw ruen'.}, welches gleichzeitig mit dem sechsten Zeitalter
(VI) stattfindet. Sowie ein \emph{8. Zeitalter}, jenes derer, die
\emph{auferstehen}\footnote{`resurgentium' und `auffersteenden'.}.

Ferner gibt es nach den Hebräern im ersten Zeitalter der Welt (I) 10
Generationen, im zweiten (II) 10, im dritten (III) 14, im vierten (IV)
17, \textbf{Matthäus} gibt jedoch um die \emph{mystische
Bedeutung}\footnote{`gratia misterij' und `aus verborgener bedeütnus',
  die historischen Aufzeichnungen hatten zwar 17 Generationen im vierten
  Zeitalter, das Matthäusevangelium (vgl. Jiménez de Cisneros,
  \href{https://doi.org/10.3931/e-rara-46695}{1517}, Matthäus 1,17)
  teilt jedoch die Genealogie Jesu absichtlich in drei Gruppen zu je 14
  ein. Dieses dient weniger zur \emph{historisch korrekten} Darstellung,
  sondern ist vielmehr eine \emph{mystische} Symbolisierung im Rahmen
  der \emph{hermeneutischen} Technik der \emph{Gematria} (wobei
  Buchstaben des Alphabets numerische Werte zugeordnet werden).
  Bemerkenswert ist hier jedoch, dass Matthäus offensichtlich mehrere
  Namen aus den historischen Aufzeichnungen \emph{weggelassen} hat, um
  die Zählung in drei Gruppen von je 14 unterteilen zu \emph{können}.}
hervorzuheben 14 an, sowie im fünften 14.

Dies sind nun die Lebensalter des Menschen: Das (I) erste ist die
\emph{frühe Kindheit}\footnote{`infantia' und `ungesprechheyt'.}, von
der Geburt bis zum 7. Lebensjahr. Das (II) zweite ist die
\emph{Kindheit}\footnote{`puericia' und `kintheit'.} bis zum 14.
Lebensjahr. Das (III) dritte ist die \emph{`Jugend'}\footnote{`adolesce{[}n{]}tia'
  und `zeittigkeit'.}, vom 15. bis zum 38. Lebensjahr. Das (IV) vierte
ist das \emph{junge Alter}\footnote{`iuue{[}n{]}t{[}us{]}' und
  `iuge{[}n{]}t'.} bis zum 49. Lebensjahr. Das (V) fünfte ist das
\emph{Alter}\footnote{`senect{[}us{]}' und `altheit'.}, vom 50. bis zum
79. Lebensjahr. Das (VI) sechste ist das \emph{hohe Alter}\footnote{`decrepita'
  und `das verlebt'.}, vom 80. Lebensjahr bis bis zu dem Ende des
Lebens.

\hypertarget{references}{%
\subsection{References}\label{references}}

Augustinus, A. (1489). \emph{De Civitate Dei}. Edited by Brant, S.,
Thomas, Nicolaus. Basel: Johannes Amerbach.
\url{https://doi.org/10.3931/e-rara-12425}

Beda. (1125). \emph{Beda, de Natura Rerum, de Temporibus, etc}. Latin
16361. Paris: Bibliothèque nationale de France. Département des
Manuscrits. \url{http://archivesetmanuscrits.bnf.fr/ark:/12148/cc76825p}

Bernard of Cluny. (1864). \emph{De Contemptu Mundi}. Harvard University:
J. T. Hayes. \url{https://books.google.com/books?id=pAM_lfiF7EUC}

De Seville, I. (1802). \emph{S. Isidori Hispalensis Episcopi Hispaniarum
Doctoris Opera Omnia}. Edited by Arevalo, F. Romae: Apvd Antonivm
Fvlgonivm. \url{https://books.google.com/books?id=wb8Z_vlSyJwC}

Dionysius. (1350). \emph{De Caelesti Hierarchia}. MS Gr 2. Oxford,
Magdalen College: Oxford Digital Library.
\url{https://digital.bodleian.ox.ac.uk/objects/ee19a692-5066-4c84-95db-d93c4fc82b97/}

Étaix, R. (1999). \emph{Gregorius Magnus, Homiliae in Evangelia}. Corpus
Christianorum Series Latina. Vol. 141. Turnhout: Brepols.
\url{https://books.google.com/books?id=r4ERnwEACAAJ}

Guazzo, F. M. (1608). \emph{Compendium Maleficarum in Tres Libros
Distinctum Ex Pluribus Authoribus}. Mediolani: Apud Haeredes August.
Tradati.
\url{https://archive.org/details/compendiummalefi00guaz/page/n2/mode/1up}

Jiménez de Cisneros, F. (1517). \emph{Biblia Polyglotta Complutensis}.
Complutum: Arnaldo Guillén de Brocar.
\url{https://doi.org/10.3931/e-rara-46695}

Meister Eckhart. (1936). \emph{Die Lateinischen Werke}. Edited by Weiss
Sturlese L. Stuttgart: Kohlhammer. \url{http://www.eckhart.de/opus.htm}

Thomas. (1485). \emph{Summa Theologica: Partes 1-3}. Basel: Michael
Wenssler. \url{https://doi.org/10.3931/e-rara-15170}

\end{document}
